In the proposed pipeline, Stage A screens epochs using a union rule over the channels supplied to the pipeline, so changing the channel set can also change which epochs reach Stage B.
Restricting the input channels therefore changes the detector itself, not merely the signals presented to a fixed detector.
The EEG channel-subset analysis accounts for this coupling by re-running the complete pipeline separately for each channel subset.
We first supply the complete 32-channel montage to the pipeline as a baseline, reported in Section~\ref{sec:exp1_reference}. This configuration assumes no prior knowledge of which brain region projects blinks most strongly, making it the most general starting point against which reduced channel sets are compared.
The analysis then evaluates anatomically defined regional subsets in Section~\ref{sec:exp1_regional} to identify where the EEG information most useful for blink detection is concentrated and to quantify the performance loss caused by excluding other scalp regions. Individual channels are subsequently assessed in Section~\ref{sec:exp1_single_channel} to determine whether a single optimally positioned electrode can preserve sufficient detection performance, thereby establishing the minimum channel configuration suitable for a simpler and more practical system.
